

\chapter{Generalised TPSAs}



\section{TPSAs}

Taylor series in one dimension:
\begin{equation}
    f(x) = \sum\limits_{k = 0}^\infty \frac{1}{k!} \, f^{(k)}(a)\,(x-a)^k\,.
\end{equation}
Taylor series in $n$ dimensions:
\begin{align*}
    f(\vec a + \Delta\vec x) & = \sum\limits_{k = 0}^\infty \frac{1}{k!} \,
    \partial_{i_1}\partial_{i_2}\cdots\partial_{i_k}\,f(\vec a)\,
    \Delta x^{i_1}\Delta x^{i_2}\cdots \Delta x^{i_k}
    \quad i_j = 1, \ldots, n\\
    & = \sum\limits_{k = 0}^\infty \frac{1}{k!} \,
    \left(\Delta x^i \partial_i\right)^{\!k} f(\vec a) \\
    & = \sum\limits_{|\alpha|\geq 0} \frac{1}{\alpha!} \,
    \Delta x^\alpha \partial^\alpha f(\vec a) \\
    \alpha & = (\alpha_1, \alpha_2, \ldots, \alpha_n) \\
    \alpha! & = \prod \alpha_j! \\
    \Delta x^\alpha & = \Delta x_1^{\alpha_1}\Delta x_2^{\alpha_2}\cdots \Delta x_n^{\alpha_n}\,.
\end{align*}

Let
\begin{equation}
    \vec r = (x,p_x,y,p_y,\zeta,\delta)^T\,.
\end{equation}
A six-dimensional Taylor map is a vector-valued map
\begin{equation}
    \mathcal M(\vec r)
    = (\mathcal M_1(\vec r),\ldots,\mathcal M_6(\vec r))^T\,,
\end{equation}
with each component represented by a TPSA.\footnote{The longitudinal canonical
coordinates are often chosen differently.}  Thus a TPSA is a scalar truncated
multivariate power series; a phase-space Taylor map is a vector of such TPSAs.

This means that each component $K$ of the map is a polynomial of order $N$:
\begin{equation}
    \mathcal{M}^K = \sum_{i + j + k + l + m + n \leq N} c_{ijklmn}\;
    x^i p_x^j y^k p_y^l \zeta^m \delta^n \,.
\end{equation}
Practically, any TPSA can be represented by a list of powers and associated coefficients for each monomial, like so:
\begin{align}
    -0.376\, xy^3 \quad\Rightarrow\quad &(1, 0, 3, 0, 0, 0):\; -0.376\,, \\
    0.8x + 0.97p_x -0.0017\, x\delta \quad\Rightarrow\quad
    &(1, 0, 0, 0, 0, 0):\; 0.8 \notag\\
    &(0, 1, 0, 0, 0, 0):\; 0.97 \notag\\
    &(1, 0, 0, 0, 0, 1):\; -0.0017\,,
\end{align}
etc. Any operation on a map is then just a matter of bookkeeping of its polynomial's powers and combining the relevant coefficients.

For instance, a derivative acts on the powers as follows:
\begin{align}\label{eq:derivatives}
    \partial_{x}\left[ (i, j, k, l, m, n):\, c \right] & = (i-1, j, k, l, m, n):\, i c \\
    \partial_{p_x}\left[ (i, j, k, l, m, n):\, c \right] & = (i, j-1, k, l, m, n):\, j c \notag\\
    \partial_{y}\left[ (i, j, k, l, m, n):\, c \right] & = (i, j, k-1, l, m, n):\, k c \notag\\
    \ldots\quad\;& \notag\\
    \partial_x^3\partial_{p_x}\left[ (i, j, k, l, m, n):\, c \right] & = (i-3, j-1, k, l, m, n):\, i(i-1)(i-2) j c \quad \text{if }i\geq 3 \text{ else }0 \notag\\
    \ldots\quad\;& \notag\,.
\end{align}



\section{Jacobian}
The Jacobian matrix of a map is given by
\begin{equation}
    \mathcal{J}_{ij} = \partial_j \mathcal{M}_i \,,
\end{equation}
which is straightforward to calculate for TPSAs using Eq.~\ref{eq:derivatives}.

The \emph{Jacobian determinant} is the determinant of the Jacobian matrix:
\begin{equation}
    \mathcal{J} = \left|\mathcal{J}_{ij}\right|\,.
\end{equation}



The Jacobian determinant alone is not a sufficient symplecticity test in more
than one degree of freedom.  For canonical coordinates, a Taylor map is
symplectic when its Jacobian matrix satisfies
\begin{equation}
    \mathcal J(\vec r)^T S\mathcal J(\vec r) = S\,,
\end{equation}
through the polynomial order available after differentiation.  In particular,
$\det\mathcal J = 1$ is a necessary consequence, but not the full condition.
